\chapter{Ansible --- Configuration Management}

\section{Presentation}

Ansible est utilise pour configurer les instances EC2 qui hebergent les conteneurs ECS.
Il intervient apres le provisionnement Terraform pour installer et configurer~:
\begin{itemize}
    \item Docker et le daemon ECS
    \item L'agent CloudWatch pour les metriques systeme
    \item Le durcissement de securite des instances
\end{itemize}

\section{Inventaire}

L'inventaire est statique (\texttt{inventory/hosts.ini}) et regroupe les
hotes en trois groupes~: \texttt{bastion}, \texttt{ecs\_instances} et
\texttt{monitoring}. Les adresses IP changent a chaque session Learner Lab
(les instances sont recreees ou redemarrees)~; l'inventaire est donc mis a
jour a partir des outputs Terraform (\texttt{bastion\_public\_ip},
IP privees des instances ECS et monitoring).

Les variables communes aux instances ECS (nom du cluster, region, log group)
sont centralisees dans \texttt{inventory/group\_vars/ecs\_instances.yml}
et doivent rester coherentes avec les outputs Terraform.

\section{Roles}

\subsection{Role \texttt{security}}

Durcissement des instances~:
\begin{itemize}
    \item Desactivation du login SSH root
    \item Desactivation de l'authentification par mot de passe
    \item Installation des mises a jour de securite automatiques
    \item Restriction des permissions cron
\end{itemize}

\subsection{Role \texttt{docker}}

Installation et activation de Docker CE sur les instances Amazon Linux 2023.

\subsection{Role \texttt{ecs-agent}}

Configuration de l'agent ECS via le template \texttt{/etc/ecs/ecs.config}
qui specifie le cluster cible et active les metadonnees de conteneur.

\subsection{Role \texttt{cloudwatch-agent}}

Installation de l'agent CloudWatch unifie avec collecte~:
\begin{itemize}
    \item Metriques CPU, memoire, disque
    \item Logs de l'agent ECS (\texttt{/var/log/ecs/ecs-agent.log})
\end{itemize}

\section{Playbook monitoring}

Un second playbook (\texttt{playbooks/monitoring.yml}) configure l'instance
dediee Prometheus/Grafana. Deux principes de securite y sont appliques~:

\begin{itemize}
    \item \textbf{Pas de secret en dur}~: le mot de passe administrateur
        Grafana et le DNS de l'ALB (cible de scrape Prometheus) sont exiges
        en \emph{extra-vars} --- une assertion en \texttt{pre\_tasks} fait
        echouer le playbook s'ils sont absents~:
\begin{lstlisting}[language=bash,caption={Invocation du playbook monitoring}]
ansible-playbook playbooks/monitoring.yml \
  -e "alb_dns_name=$(terraform -chdir=../terraform/environments/dev \
      output -raw alb_dns_name)" \
  -e "grafana_admin_password=$GRAFANA_ADMIN_PASSWORD"
\end{lstlisting}
    \item \textbf{Versions epinglees}~: les images Docker
        (\texttt{prom/prometheus:v2.53.4}, \texttt{grafana/grafana:11.6.1},
        \texttt{prom/node-exporter:v1.8.2}) sont definies dans les
        \texttt{defaults} de chaque role plutot qu'en \texttt{latest},
        garantissant des deploiements reproductibles. Un changement de
        version detecte a l'execution recree le conteneur concerne.
\end{itemize}

\section{Connexion via Bastion}

Les instances ECS sont dans des sous-reseaux prives. Ansible s'y connecte
via SSH ProxyJump a travers le bastion~:

\begin{lstlisting}[language=bash,caption={Connexion SSH via Bastion}]
ansible_ssh_common_args: >
  -o ProxyJump=ec2-user@<bastion_ip>
\end{lstlisting}

\begin{figure}[H]
    \centering
    \includegraphics[width=\textwidth]{ansible_init.png}
    \caption{Execution du playbook Ansible}
\end{figure}
